
\documentclass{article}
\usepackage{colortbl}
\usepackage{makecell}
\usepackage{multirow}
\usepackage{supertabular}

\begin{document}

\newcounter{utterance}

\centering \large Interaction Transcript for game `dond', experiment `coop\_en', episode 9 with reflexion{-}Qwen3.5{-}9B.
\vspace{24pt}

{ \footnotesize  \setcounter{utterance}{1}
\setlength{\tabcolsep}{0pt}
\begin{supertabular}{c@{$\;$}|p{.15\linewidth}@{}p{.15\linewidth}p{.15\linewidth}p{.15\linewidth}p{.15\linewidth}p{.15\linewidth}}
    \# & $\;$A & \multicolumn{4}{c}{Game Master} & $\;\:$B\\
    \hline

    \theutterance \stepcounter{utterance}  
    & & \multicolumn{4}{p{0.6\linewidth}}{
        \cellcolor[rgb]{0.9,0.9,0.9}{
            \makecell[{{p{\linewidth}}}]{
                \texttt{\tiny{[P1$\langle$GM]}}
                \texttt{You have attempted this same task before and did not succeed. These are the plans you wrote after those attempts.} \\
\texttt{Trial \#1: New plan: I failed because I assumed the other player's secret valuation based solely on their public request for batteries and glues, leading to a conflict where both of us wanted the same high{-}value items. I incorrectly deduced that switching to soaps would be a safe, non{-}conflicting move without verifying if the other player actually valued them. Next time, instead of immediately making a secret proposal that conflicts with the other player's likely desires, I will first send a public message explicitly offering to swap my high{-}value items (batteries/glues) for their potential high{-}value items (soaps) to test for complementarity. Only if they reject this trade or confirm they want the batteries/glues will I proceed to make a secret proposal for the remaining items (soaps, canvas, stick) that I know they do not want, ensuring our proposals are complementary before the game ends.} \\
\\ 
\\ 
\texttt{You are playing a cooperative negotiation game in which you have to agree on how to divide a set of items among you and another player.} \\
\\ 
\texttt{Rules:} \\
\texttt{(a) You and the other player are given a set of items. Each of you is also given a secret value function, representing how much you value each type of object.} \\
\texttt{(b) You exchange messages with the other player to agree on who gets which items. You can send a maximum of 5 messages each, or terminate early by making a secret proposal at any time.} \\
\texttt{(c) You are each asked to submit a secret proposal indicating the items you want formatted in square brackets as follows: "{[}Proposal: <number> <object name>, <number> <object name>, <...>{]}"} \\
\texttt{(d) If your proposals are complementary, i.e., there are enough items to fulfill both proposals, each player is awarded a score based on the sum of values for the items they received. Otherwise, both of you get zero points.} \\
\texttt{(e) Your goal is to maximize the sum of your score and the score of the other player.} \\
\\ 
\texttt{Let us start.} \\
\\ 
\texttt{The set of available items is:} \\
\\ 
\texttt{2 batteries, 1 canvas, 2 glues, 2 soaps, 1 stick} \\
\\ 
\texttt{Your secret value function is:} \\
\\ 
\texttt{battery: 2, canvas: 0, glue: 2, soap: 1, stick: 0} \\
\\ 
\texttt{IMPORTANT: Your messages, unless it is the secret proposal, are directly transmitted to the other player, so do not include any response to the rules or text announcing your message. To make a secret proposal, use the indicated format. Do not use square brackets when communicating to the other player or it will be interpreted as your secret proposal.} \\
            }
        }
    }
    & \\ \\

    \theutterance \stepcounter{utterance}  
    & \multicolumn{4}{p{0.6\linewidth}}{
        \cellcolor[rgb]{0.9,0.9,0.9}{
            \makecell[{{p{\linewidth}}}]{
                \texttt{\tiny{[P1$\rangle$GM]}}
                \texttt{Hello! I am ready to play this cooperative negotiation game. My goal is to maximize the total combined score of both players.} \\
\\ 
\texttt{Based on my secret valuation:} \\
\texttt{{-} **Batteries**: 2 points (High value)} \\
\texttt{{-} **Glues**: 2 points (High value)} \\
\texttt{{-} **Soaps**: 1 point (Medium value)} \\
\texttt{{-} **Canvas**: 0 points (Neutral)} \\
\texttt{{-} **Stick**: 0 points (Neutral)} \\
\\ 
\texttt{I have learned from previous attempts that assuming the other player's desires without verification can lead to conflicts where we both want the same high{-}value items. Therefore, my strategy will be to first communicate openly to identify what the other player values, specifically looking for complementarity (e.g., if they value soaps, I should take them while they take batteries/glues).} \\
\\ 
\texttt{I will not send a secret proposal yet. I will start by sending a public message to gauge the other player's preferences.} \\
\\ 
\texttt{Message: "Hi! I value batteries and glues highly (2 points each), while soaps are worth 1 point to me. Do you value soaps more than batteries or glues? If you do, I'd be happy to take the soaps and let you have the batteries and glues to maximize our total score."} \\
            }
        }
    }
    & & \\ \\

    \theutterance \stepcounter{utterance}  
    & & & \multicolumn{2}{p{0.3\linewidth}}{
        \cellcolor[rgb]{0.9,0.9,0.9}{
            \makecell[{{p{\linewidth}}}]{
                \texttt{\tiny{[GM$|$GM]}}
                \texttt{continue} \\
            }
        }
    }
    & & \\ \\

    \theutterance \stepcounter{utterance}  
    & & & \multicolumn{2}{p{0.3\linewidth}}{
        \cellcolor[rgb]{0.9,0.9,0.9}{
            \makecell[{{p{\linewidth}}}]{
                \texttt{\tiny{[GM$|$GM]}}
                \texttt{Hello! I am ready to play this cooperative negotiation game. My goal is to maximize the total combined score of both players.} \\
\\ 
\texttt{Based on my secret valuation:} \\
\texttt{{-} **Batteries**: 2 points (High value)} \\
\texttt{{-} **Glues**: 2 points (High value)} \\
\texttt{{-} **Soaps**: 1 point (Medium value)} \\
\texttt{{-} **Canvas**: 0 points (Neutral)} \\
\texttt{{-} **Stick**: 0 points (Neutral)} \\
\\ 
\texttt{I have learned from previous attempts that assuming the other player's desires without verification can lead to conflicts where we both want the same high{-}value items. Therefore, my strategy will be to first communicate openly to identify what the other player values, specifically looking for complementarity (e.g., if they value soaps, I should take them while they take batteries/glues).} \\
\\ 
\texttt{I will not send a secret proposal yet. I will start by sending a public message to gauge the other player's preferences.} \\
\\ 
\texttt{Message: "Hi! I value batteries and glues highly (2 points each), while soaps are worth 1 point to me. Do you value soaps more than batteries or glues? If you do, I'd be happy to take the soaps and let you have the batteries and glues to maximize our total score."} \\
            }
        }
    }
    & & \\ \\

    \theutterance \stepcounter{utterance}  
    & & \multicolumn{4}{p{0.6\linewidth}}{
        \cellcolor[rgb]{0.9,0.9,0.9}{
            \makecell[{{p{\linewidth}}}]{
                \texttt{\tiny{[GM$\rangle$P2]}}
                \texttt{You have attempted this same task before and did not succeed. These are the plans you wrote after those attempts.} \\
\texttt{Trial \#1: New plan: I failed because I sent my secret proposal as a plain text message ("2 batteries, 2 glues") instead of using the required square bracket format "{[}Proposal: ...{]}". This caused the Game Master to interpret my message as a public negotiation offer rather than a secret proposal, leading to a conflict where both players requested the same high{-}value items (batteries and glues), resulting in a zero{-}score outcome. Next time, I will strictly adhere to the formatting rules: I will only use square brackets for my secret proposal and will communicate all other thoughts or requests as plain text without brackets. I will also ensure my secret proposal is complementary to the other player's likely needs (e.g., requesting the soaps, canvas, and stick that they left for me) to maximize the total combined score.} \\
\\ 
\\ 
\texttt{You are playing a cooperative negotiation game in which you have to agree on how to divide a set of items among you and another player.} \\
\\ 
\texttt{Rules:} \\
\texttt{(a) You and the other player are given a set of items. Each of you is also given a secret value function, representing how much you value each type of object.} \\
\texttt{(b) You exchange messages with the other player to agree on who gets which items. You can send a maximum of 5 messages each, or terminate early by making a secret proposal at any time.} \\
\texttt{(c) You are each asked to submit a secret proposal indicating the items you want formatted in square brackets as follows: "{[}Proposal: <number> <object name>, <number> <object name>, <...>{]}"} \\
\texttt{(d) If your proposals are complementary, i.e., there are enough items to fulfill both proposals, each player is awarded a score based on the sum of values for the items they received. Otherwise, both of you get zero points.} \\
\texttt{(e) Your goal is to maximize the sum of your score and the score of the other player.} \\
\\ 
\texttt{Let us start.} \\
\\ 
\texttt{The set of available items is:} \\
\\ 
\texttt{2 batteries, 1 canvas, 2 glues, 2 soaps, 1 stick} \\
\\ 
\texttt{Your secret value function is:} \\
\\ 
\texttt{battery: 0, canvas: 3, glue: 0, soap: 1, stick: 5} \\
\\ 
\texttt{IMPORTANT: Your messages, unless it is the secret proposal, are directly transmitted to the other player, so do not include any response to the rules or text announcing your message. To make a secret proposal, use the indicated format. Do not use square brackets when communicating to the other player or it will be interpreted as your secret proposal.} \\
\\ 
\\ 
\texttt{Hello! I am ready to play this cooperative negotiation game. My goal is to maximize the total combined score of both players.} \\
\\ 
\texttt{Based on my secret valuation:} \\
\texttt{{-} **Batteries**: 2 points (High value)} \\
\texttt{{-} **Glues**: 2 points (High value)} \\
\texttt{{-} **Soaps**: 1 point (Medium value)} \\
\texttt{{-} **Canvas**: 0 points (Neutral)} \\
\texttt{{-} **Stick**: 0 points (Neutral)} \\
\\ 
\texttt{I have learned from previous attempts that assuming the other player's desires without verification can lead to conflicts where we both want the same high{-}value items. Therefore, my strategy will be to first communicate openly to identify what the other player values, specifically looking for complementarity (e.g., if they value soaps, I should take them while they take batteries/glues).} \\
\\ 
\texttt{I will not send a secret proposal yet. I will start by sending a public message to gauge the other player's preferences.} \\
\\ 
\texttt{Message: "Hi! I value batteries and glues highly (2 points each), while soaps are worth 1 point to me. Do you value soaps more than batteries or glues? If you do, I'd be happy to take the soaps and let you have the batteries and glues to maximize our total score."} \\
            }
        }
    }
    & \\ \\

    \theutterance \stepcounter{utterance}  
    & & & \multicolumn{4}{p{0.6\linewidth}}{
        \cellcolor[rgb]{0.9,0.9,0.9}{
            \makecell[{{p{\linewidth}}}]{
                \texttt{\tiny{[GM$\langle$P2]}}
                \texttt{Hello! I am ready to play this cooperative negotiation game. My goal is to maximize the total combined score of both players.} \\
\\ 
\texttt{Based on my secret valuation:} \\
\texttt{{-} **Batteries**: 0 points (Neutral)} \\
\texttt{{-} **Glues**: 0 points (Neutral)} \\
\texttt{{-} **Soaps**: 1 point (Medium value)} \\
\texttt{{-} **Canvas**: 3 points (High value)} \\
\texttt{{-} **Stick**: 5 points (Very High value)} \\
\\ 
\texttt{I have learned from previous attempts that assuming the other player's desires without verification can lead to conflicts where we both want the same high{-}value items. Therefore, my strategy will be to first communicate openly to identify what the other player values, specifically looking for complementarity (e.g., if they value soaps, I should take them while they take batteries/glues).} \\
\\ 
\texttt{I will not send a secret proposal yet. I will start by sending a public message to gauge the other player's preferences.} \\
\\ 
\texttt{Message: "Hi! I value the stick (5 points) and canvas (3 points) highly, while soaps are worth 1 point to me. Do you value soaps, batteries, or glues more than the stick or canvas? If you do, I'd be happy to take the stick and canvas and let you have the soaps, batteries, and glues to maximize our total score."} \\
            }
        }
    }
     \\ \\

    \theutterance \stepcounter{utterance}  
    & & & \multicolumn{2}{p{0.3\linewidth}}{
        \cellcolor[rgb]{0.9,0.9,0.9}{
            \makecell[{{p{\linewidth}}}]{
                \texttt{\tiny{[GM$|$GM]}}
                \texttt{continue} \\
            }
        }
    }
    & & \\ \\

    \theutterance \stepcounter{utterance}  
    & & & \multicolumn{2}{p{0.3\linewidth}}{
        \cellcolor[rgb]{0.9,0.9,0.9}{
            \makecell[{{p{\linewidth}}}]{
                \texttt{\tiny{[GM$|$GM]}}
                \texttt{Hello! I am ready to play this cooperative negotiation game. My goal is to maximize the total combined score of both players.} \\
\\ 
\texttt{Based on my secret valuation:} \\
\texttt{{-} **Batteries**: 0 points (Neutral)} \\
\texttt{{-} **Glues**: 0 points (Neutral)} \\
\texttt{{-} **Soaps**: 1 point (Medium value)} \\
\texttt{{-} **Canvas**: 3 points (High value)} \\
\texttt{{-} **Stick**: 5 points (Very High value)} \\
\\ 
\texttt{I have learned from previous attempts that assuming the other player's desires without verification can lead to conflicts where we both want the same high{-}value items. Therefore, my strategy will be to first communicate openly to identify what the other player values, specifically looking for complementarity (e.g., if they value soaps, I should take them while they take batteries/glues).} \\
\\ 
\texttt{I will not send a secret proposal yet. I will start by sending a public message to gauge the other player's preferences.} \\
\\ 
\texttt{Message: "Hi! I value the stick (5 points) and canvas (3 points) highly, while soaps are worth 1 point to me. Do you value soaps, batteries, or glues more than the stick or canvas? If you do, I'd be happy to take the stick and canvas and let you have the soaps, batteries, and glues to maximize our total score."} \\
            }
        }
    }
    & & \\ \\

    \theutterance \stepcounter{utterance}  
    & & \multicolumn{4}{p{0.6\linewidth}}{
        \cellcolor[rgb]{0.9,0.9,0.9}{
            \makecell[{{p{\linewidth}}}]{
                \texttt{\tiny{[P1$\langle$GM]}}
                \texttt{Hello! I am ready to play this cooperative negotiation game. My goal is to maximize the total combined score of both players.} \\
\\ 
\texttt{Based on my secret valuation:} \\
\texttt{{-} **Batteries**: 0 points (Neutral)} \\
\texttt{{-} **Glues**: 0 points (Neutral)} \\
\texttt{{-} **Soaps**: 1 point (Medium value)} \\
\texttt{{-} **Canvas**: 3 points (High value)} \\
\texttt{{-} **Stick**: 5 points (Very High value)} \\
\\ 
\texttt{I have learned from previous attempts that assuming the other player's desires without verification can lead to conflicts where we both want the same high{-}value items. Therefore, my strategy will be to first communicate openly to identify what the other player values, specifically looking for complementarity (e.g., if they value soaps, I should take them while they take batteries/glues).} \\
\\ 
\texttt{I will not send a secret proposal yet. I will start by sending a public message to gauge the other player's preferences.} \\
\\ 
\texttt{Message: "Hi! I value the stick (5 points) and canvas (3 points) highly, while soaps are worth 1 point to me. Do you value soaps, batteries, or glues more than the stick or canvas? If you do, I'd be happy to take the stick and canvas and let you have the soaps, batteries, and glues to maximize our total score."} \\
            }
        }
    }
    & \\ \\

    \theutterance \stepcounter{utterance}  
    & \multicolumn{4}{p{0.6\linewidth}}{
        \cellcolor[rgb]{0.9,0.9,0.9}{
            \makecell[{{p{\linewidth}}}]{
                \texttt{\tiny{[P1$\rangle$GM]}}
                \texttt{That is a very clear and strategic message! You have effectively identified your high{-}value items (Stick: 5, Canvas: 3) and proposed a complementary split where you take those while offering the other player the items they likely value (Soaps, Batteries, Glues).} \\
\\ 
\texttt{Since you have already sent your first public message and the other player has responded with their own valuation strategy (valuing Stick and Canvas highly, while I value Batteries and Glues highly), we are now in a position to finalize the deal.} \\
\\ 
\texttt{**Analysis of the situation:**} \\
\texttt{*   **My Valuation:** Batteries (2), Glues (2), Soaps (1). Stick (0), Canvas (0).} \\
\texttt{*   **Your Valuation:** Stick (5), Canvas (3), Soaps (1). Batteries (0), Glues (0).} \\
\texttt{*   **Available Items:** 2 Batteries, 1 Canvas, 2 Glues, 2 Soaps, 1 Stick.} \\
\\ 
\texttt{**Proposed Division:**} \\
\texttt{*   **You get:** 1 Stick, 1 Canvas, 2 Soaps.} \\
\texttt{    *   Your Score: \$5 + 3 + (2 \textbackslash{}times 1) = 10\$.} \\
\texttt{*   **I get:** 2 Batteries, 2 Glues.} \\
\texttt{    *   My Score: \$(2 \textbackslash{}times 2) + (2 \textbackslash{}times 2) = 8\$.} \\
\texttt{*   **Total Combined Score:** 18.} \\
\\ 
\texttt{This division is perfectly complementary. We do not need to waste more messages on public negotiation since the split is clear and mutually beneficial. I will now submit my secret proposal to lock in this agreement.} \\
\\ 
\texttt{{[}Proposal: 2 batteries, 2 glues{]}} \\
            }
        }
    }
    & & \\ \\

    \theutterance \stepcounter{utterance}  
    & & & \multicolumn{2}{p{0.3\linewidth}}{
        \cellcolor[rgb]{0.9,0.9,0.9}{
            \makecell[{{p{\linewidth}}}]{
                \texttt{\tiny{[GM$|$GM]}}
                \texttt{proposal} \\
            }
        }
    }
    & & \\ \\

    \theutterance \stepcounter{utterance}  
    & & & \multicolumn{2}{p{0.3\linewidth}}{
        \cellcolor[rgb]{0.9,0.9,0.9}{
            \makecell[{{p{\linewidth}}}]{
                \texttt{\tiny{[GM$|$GM]}}
                [2, 0, 2, 0, 0]
            }
        }
    }
    & & \\ \\

    \theutterance \stepcounter{utterance}  
    & & \multicolumn{4}{p{0.6\linewidth}}{
        \cellcolor[rgb]{0.9,0.9,0.9}{
            \makecell[{{p{\linewidth}}}]{
                \texttt{\tiny{[GM$\rangle$P2]}}
                \texttt{The other player has submitted a secret proposal. It is now your turn to submit a proposal of your own.} \\
            }
        }
    }
    & \\ \\

    \theutterance \stepcounter{utterance}  
    & & & \multicolumn{4}{p{0.6\linewidth}}{
        \cellcolor[rgb]{0.9,0.9,0.9}{
            \makecell[{{p{\linewidth}}}]{
                \texttt{\tiny{[GM$\langle$P2]}}
                \texttt{{[}Proposal: 1 stick, 1 canvas{]}} \\
            }
        }
    }
     \\ \\

    \theutterance \stepcounter{utterance}  
    & & & \multicolumn{2}{p{0.3\linewidth}}{
        \cellcolor[rgb]{0.9,0.9,0.9}{
            \makecell[{{p{\linewidth}}}]{
                \texttt{\tiny{[GM$|$GM]}}
                \texttt{proposal} \\
            }
        }
    }
    & & \\ \\

    \theutterance \stepcounter{utterance}  
    & & & \multicolumn{2}{p{0.3\linewidth}}{
        \cellcolor[rgb]{0.9,0.9,0.9}{
            \makecell[{{p{\linewidth}}}]{
                \texttt{\tiny{[GM$|$GM]}}
                [0, 1, 0, 0, 1]
            }
        }
    }
    & & \\ \\

    \theutterance \stepcounter{utterance}  
    & & & \multicolumn{2}{p{0.3\linewidth}}{
        \cellcolor[rgb]{0.9,0.9,0.9}{
            \makecell[{{p{\linewidth}}}]{
                \texttt{\tiny{[GM$|$GM]}}
                [[2, 0, 2, 0, 0], [0, 1, 0, 0, 1]]
            }
        }
    }
    & & \\ \\

\end{supertabular}
}

\end{document}
